% 图契约：本题的源、边界和推导逐项对应；箭头符号按图注解释。
\documentclass[tikz,border=7pt]{standalone}
\usepackage{fontspec,xeCJK,amsmath,amssymb}
\IfFontExistsTF{Times New Roman}{\setmainfont{Times New Roman}}{\setmainfont{TeX Gyre Termes}}
\IfFontExistsTF{Heiti SC}{\setCJKmainfont[AutoFakeBold=1.5]{Heiti SC}}{\setCJKmainfont[AutoFakeBold=1.5]{FandolHei-Regular}}
\usetikzlibrary{arrows.meta,bending,calc,patterns,positioning}
\definecolor{ink}{HTML}{243745}
\definecolor{blue}{HTML}{275DAB}
\definecolor{green}{HTML}{177C63}
\definecolor{red}{HTML}{BD4935}
\definecolor{orange}{HTML}{AA6C16}
\definecolor{purple}{HTML}{715393}
\tikzset{every node/.style={font=\fontsize{9.3}{12.5}\selectfont,text=ink},
axis/.style={-{Stealth[length=4pt]},line width=.65pt,draw=ink!65},
vec/.style={-{Stealth[length=5pt]},line width=1pt,draw=blue},
short/.style={-{Stealth[length=3pt]},line width=.8pt},
guide/.style={dashed,draw=ink!40,line width=.5pt},
card/.style={rounded corners=2pt,draw=blue!65,fill=blue!4,align=center,text width=6.5cm,inner xsep=9pt,inner ysep=7pt},
title/.style={font=\fontsize{10.5}{14}\selectfont\bfseries,text width=15cm,align=center},
label/.style={font=\fontsize{8.4}{11}\selectfont,align=center},
note/.style={label,text width=6.2cm}}
\begin{document}
\begin{tikzpicture}[x=1cm,y=1cm]
\path[use as bounding box] (-8,-3.5) rectangle(8,3.5);
\node[title] at (0,3.15) {M3-014a\quad 错开带电球得到均匀退极化场};
\begin{scope}[shift={(-4,0)}]

\draw[blue,thick] (-.25,0)circle(1.55);\draw[red,thick] (.25,0)circle(1.55);
\fill[blue](-.25,0)circle(2pt) node[below left] {$O_-$};\fill[red](.25,0)circle(2pt) node[below right] {$O_+$};
\fill[ink](.55,.85)circle(2pt) node[above] {$A$};
\draw[guide] (-.25,0)--(.55,.85) node[midway,left] {$\mathbf r$};
\draw[guide] (.25,0)--(.55,.85) node[midway,right] {$\mathbf r-\mathbf d$};
\draw[vec,green] (-.25,-1.9)--(.25,-1.9) node[right] {$\mathbf d$};
\node[label] at(-1.4,1.8){$-\rho$};\node[label] at(1.4,1.8){$+\rho$};
\node[note] at(0,-2.6){微小错位已夸大；取保持 $\rho\mathbf d=\mathbf P$ 的极限};
\end{scope}
\node[card,draw=blue!65,fill=blue!4] at (3.9,1.8) {$\mathbf E_+=\rho(\mathbf r-\mathbf d)/(3\varepsilon_0)$};
\node[card,draw=green!65,fill=green!4] at (3.9,0) {$\mathbf E_-=-\rho\mathbf r/(3\varepsilon_0)$};
\node[card,draw=orange!65,fill=orange!4] at (3.9,-1.8) {$\mathbf E_{\rm dep}=-\rho\mathbf d/(3\varepsilon_0)$\\$=-\mathbf P/(3\varepsilon_0)$};
\end{tikzpicture}\end{document}